\subsection{Xác định thông tin cơ bản cho nghiệp vụ của bài toán}

Hệ thống quản lý phòng khám nha khoa được thiết kế cho các phòng khám nhỏ với hai nhóm người dùng chính: Bác sĩ (Owner/Doctor) và Nhân viên (Staff). Dưới đây là phân tích các nghiệp vụ chính của hệ thống:

\begin{table}[h]
\centering
\caption{Phân tích nghiệp vụ quản lý bệnh nhân}
\begin{tabular}{|p{4cm}|p{5cm}|p{4.5cm}|}
\hline
\textbf{Input} & \textbf{Process} & \textbf{Output} \\
\hline
Thông tin bệnh nhân (họ tên, số điện thoại, địa chỉ, ngày sinh, ghi chú) & Đăng ký và lưu trữ thông tin bệnh nhân vào cơ sở dữ liệu, kiểm tra tính hợp lệ của dữ liệu & Hồ sơ bệnh nhân được tạo, mã ID bệnh nhân, thông báo thành công \\
\hline
Mã ID bệnh nhân, từ khóa tìm kiếm & Tìm kiếm bệnh nhân theo tên, số điện thoại hoặc mã ID & Danh sách bệnh nhân phù hợp, thông tin chi tiết bệnh nhân \\
\hline
Mã ID bệnh nhân, thông tin cập nhật & Cập nhật thông tin bệnh nhân trong hệ thống & Hồ sơ bệnh nhân được cập nhật, thông báo thành công \\
\hline
Mã ID bệnh nhân & Xóa thông tin bệnh nhân khỏi hệ thống & Xác nhận xóa, cập nhật cơ sở dữ liệu \\
\hline
\end{tabular}
\end{table}

\begin{table}[h]
\centering
\caption{Phân tích nghiệp vụ quản lý lịch hẹn}
\begin{tabular}{|p{4cm}|p{5cm}|p{4.5cm}|}
\hline
\textbf{Input} & \textbf{Process} & \textbf{Output} \\
\hline
ID bệnh nhân, ID bác sĩ, thời gian hẹn, mô tả & Tạo lịch hẹn mới, kiểm tra xung đột lịch, gán trạng thái "scheduled" & Lịch hẹn được tạo với mã ID, thông báo xác nhận \\
\hline
Ngày, tuần, tháng cần xem & Truy vấn lịch hẹn theo khoảng thời gian, sắp xếp theo thời gian & Danh sách lịch hẹn, hiển thị dạng lịch hoặc danh sách \\
\hline
Mã ID lịch hẹn, trạng thái mới & Cập nhật trạng thái lịch hẹn (scheduled, completed, cancelled) & Lịch hẹn được cập nhật, thông báo cho các bên liên quan \\
\hline
\end{tabular}
\end{table}

\begin{table}[h]
\centering
\caption{Phân tích nghiệp vụ quản lý điều trị}
\begin{tabular}{|p{4cm}|p{5cm}|p{4.5cm}|}
\hline
\textbf{Input} & \textbf{Process} & \textbf{Output} \\
\hline
ID bệnh nhân, ID bác sĩ, ngày điều trị, mô tả, tổng chi phí & Tạo hồ sơ điều trị, liên kết với bệnh nhân và bác sĩ & Hồ sơ điều trị được tạo, mã ID điều trị \\
\hline
Mã ID điều trị & Xem chi tiết điều trị, lịch sử thanh toán & Thông tin điều trị chi tiết, danh sách thanh toán liên quan \\
\hline
Mã ID điều trị, số tiền, phương thức thanh toán & Ghi nhận thanh toán, cập nhật công nợ & Phiếu thu, cập nhật trạng thái thanh toán \\
\hline
Mã ID điều trị, thông tin cập nhật & Sửa đổi thông tin điều trị & Hồ sơ điều trị được cập nhật \\
\hline
\end{tabular}
\end{table}

\begin{table}[h]
\centering
\caption{Phân tích nghiệp vụ quản lý tồn kho}
\begin{tabular}{|p{4cm}|p{5cm}|p{4.5cm}|}
\hline
\textbf{Input} & \textbf{Process} & \textbf{Output} \\
\hline
Thông tin vật tư (tên, danh mục, số lượng, giá, ngày hết hạn) & Thêm vật tư vào kho, phân loại theo danh mục (thuốc, vật tư, thiết bị) & Mặt hàng tồn kho được tạo, cảnh báo mức tồn kho tối thiểu \\
\hline
Mã ID vật tư, số lượng nhập/xuất & Ghi nhận giao dịch nhập/xuất kho, cập nhật số lượng tồn & Phiếu nhập/xuất kho, cập nhật số lượng tồn kho \\
\hline
Ngày, khoảng thời gian & Kiểm tra vật tư sắp hết hạn, tồn kho dưới mức tối thiểu & Danh sách cảnh báo, báo cáo tồn kho \\
\hline
Thông tin nhà cung cấp & Quản lý danh sách nhà cung cấp & Hồ sơ nhà cung cấp, thông tin liên hệ \\
\hline
\end{tabular}
\end{table}

\begin{table}[h]
\centering
\caption{Phân tích nghiệp vụ quản lý đơn đặt labo}
\begin{tabular}{|p{4cm}|p{5cm}|p{4.5cm}|}
\hline
\textbf{Input} & \textbf{Process} & \textbf{Output} \\
\hline
ID điều trị, ID đối tác labo, mô tả, giá, ngày giao & Tạo đơn đặt labo, gán trạng thái "ORDERED" & Đơn đặt labo với mã ID, xác nhận đặt hàng \\
\hline
Mã ID đơn labo & Xem chi tiết đơn labo, theo dõi trạng thái & Thông tin đơn labo, trạng thái hiện tại \\
\hline
Mã ID đơn labo, trạng thái mới & Cập nhật trạng thái (ORDERED, RECEIVED, INSTALLED, CANCELED) & Đơn labo được cập nhật, thông báo cho bác sĩ \\
\hline
\end{tabular}
\end{table}

\begin{table}[h]
\centering
\caption{Phân tích nghiệp vụ báo cáo thống kê}
\begin{tabular}{|p{4cm}|p{5cm}|p{4.5cm}|}
\hline
\textbf{Input} & \textbf{Process} & \textbf{Output} \\
\hline
Khoảng thời gian (tuần, tháng, năm) & Tính toán doanh thu, chi phí, lợi nhuận theo ngày & Báo cáo doanh thu chi tiết, biểu đồ doanh thu theo ngày \\
\hline
Khoảng thời gian, ID nhân viên & Phân tích hiệu suất làm việc của nhân viên & Báo cáo số lượng điều trị, doanh thu theo nhân viên \\
\hline
ID phòng khám, khoảng thời gian & Thống kê thanh toán, công nợ & Báo cáo tài chính tổng hợp \\
\hline
\end{tabular}
\end{table}

\subsection{Xây dựng biểu đồ mô tả nghiệp vụ và phân cấp chức năng}

\textbf{Biểu đồ hoạt động mô tả nghiệp vụ chính:}

Do hạn chế về công cụ sinh biểu đồ, dưới đây là mô tả luồng hoạt động chính của hệ thống dưới dạng văn bản có cấu trúc:

\textbf{Quy trình điều trị bệnh nhân (Treatment Workflow):}
\begin{enumerate}
\item \textbf{Bắt đầu:} Bệnh nhân đến phòng khám
\item \textbf{Kiểm tra:} Bệnh nhân đã có hồ sơ?
    \begin{itemize}
    \item Nếu KHÔNG: Nhân viên đăng ký hồ sơ bệnh nhân mới $\rightarrow$ Lưu vào hệ thống
    \item Nếu CÓ: Tìm kiếm hồ sơ bệnh nhân trong hệ thống
    \end{itemize}
\item \textbf{Kiểm tra:} Có lịch hẹn trước không?
    \begin{itemize}
    \item Nếu CÓ: Cập nhật trạng thái lịch hẹn thành "completed"
    \item Nếu KHÔNG: Tiếp tục quy trình
    \end{itemize}
\item \textbf{Hoạt động:} Bác sĩ khám và điều trị bệnh nhân
\item \textbf{Hoạt động:} Tạo hồ sơ điều trị (Treatment Record)
    \begin{itemize}
    \item Nhập thông tin: Mô tả triệu chứng, chẩn đoán, phương pháp điều trị
    \item Nhập tổng chi phí điều trị
    \end{itemize}
\item \textbf{Quyết định:} Cần đặt hàng labo?
    \begin{itemize}
    \item Nếu CÓ: Tạo đơn đặt labo $\rightarrow$ Gửi cho đối tác labo
    \item Nếu KHÔNG: Tiếp tục
    \end{itemize}
\item \textbf{Quyết định:} Cần sử dụng vật tư từ kho?
    \begin{itemize}
    \item Nếu CÓ: Tạo phiếu xuất kho $\rightarrow$ Cập nhật tồn kho
    \item Nếu KHÔNG: Tiếp tục
    \end{itemize}
\item \textbf{Hoạt động:} Thanh toán
    \begin{itemize}
    \item Bệnh nhân thanh toán toàn bộ $\rightarrow$ Tạo phiếu thu
    \item Bệnh nhân thanh toán một phần $\rightarrow$ Ghi nhận công nợ
    \end{itemize}
\item \textbf{Quyết định:} Cần hẹn tái khám?
    \begin{itemize}
    \item Nếu CÓ: Tạo lịch hẹn cho lần khám tiếp theo
    \item Nếu KHÔNG: Tiếp tục
    \end{itemize}
\item \textbf{Kết thúc:} Hoàn thành quy trình điều trị
\end{enumerate}

\vspace{1em}

\textbf{Biểu đồ phân cấp chức năng (Business Function Diagram):}

Hệ thống được tổ chức theo cấu trúc phân cấp như sau:

\textbf{Cấp 0: Hệ thống quản lý phòng khám nha khoa}

\textbf{Cấp 1: Các chức năng chính}
\begin{itemize}
\item 1. Quản lý phòng khám
\item 2. Quản lý bệnh nhân
\item 3. Quản lý lịch hẹn
\item 4. Quản lý điều trị
\item 5. Quản lý tồn kho
\item 6. Quản lý đơn đặt labo
\item 7. Báo cáo và thống kê
\item 8. Quản lý người dùng
\end{itemize}

\textbf{Cấp 2: Phân rã chức năng chi tiết}

\textbf{1. Quản lý phòng khám}
\begin{itemize}
\item 1.1. Tạo phòng khám mới
\item 1.2. Tham gia phòng khám (bằng mã code)
\item 1.3. Quản lý thành viên
    \begin{itemize}
    \item 1.3.1. Duyệt yêu cầu tham gia
    \item 1.3.2. Thiết lập lương cho thành viên
    \item 1.3.3. Xóa thành viên
    \end{itemize}
\item 1.4. Cập nhật thông tin phòng khám
\end{itemize}

\textbf{2. Quản lý bệnh nhân}
\begin{itemize}
\item 2.1. Thêm bệnh nhân mới
\item 2.2. Tìm kiếm bệnh nhân
\item 2.3. Xem chi tiết hồ sơ bệnh nhân
\item 2.4. Cập nhật thông tin bệnh nhân
\item 2.5. Xóa bệnh nhân
\item 2.6. Xem lịch sử điều trị
\end{itemize}

\textbf{3. Quản lý lịch hẹn}
\begin{itemize}
\item 3.1. Tạo lịch hẹn mới
\item 3.2. Xem lịch hẹn (dạng danh sách)
\item 3.3. Xem lịch hẹn (dạng lịch)
\item 3.4. Cập nhật trạng thái lịch hẹn
\item 3.5. Hủy lịch hẹn
\item 3.6. Tìm kiếm lịch hẹn
\end{itemize}

\textbf{4. Quản lý điều trị}
\begin{itemize}
\item 4.1. Tạo hồ sơ điều trị
\item 4.2. Xem danh sách điều trị
\item 4.3. Xem chi tiết điều trị
\item 4.4. Cập nhật thông tin điều trị
\item 4.5. Quản lý thanh toán
    \begin{itemize}
    \item 4.5.1. Ghi nhận thanh toán
    \item 4.5.2. Xem lịch sử thanh toán
    \item 4.5.3. Theo dõi công nợ
    \end{itemize}
\end{itemize}

\textbf{5. Quản lý tồn kho}
\begin{itemize}
\item 5.1. Quản lý vật tư
    \begin{itemize}
    \item 5.1.1. Thêm vật tư mới
    \item 5.1.2. Cập nhật thông tin vật tư
    \item 5.1.3. Xóa vật tư
    \end{itemize}
\item 5.2. Quản lý giao dịch kho
    \begin{itemize}
    \item 5.2.1. Nhập kho
    \item 5.2.2. Xuất kho
    \item 5.2.3. Xem lịch sử giao dịch
    \end{itemize}
\item 5.3. Quản lý nhà cung cấp
    \begin{itemize}
    \item 5.3.1. Thêm nhà cung cấp
    \item 5.3.2. Cập nhật thông tin
    \item 5.3.3. Xóa nhà cung cấp
    \end{itemize}
\item 5.4. Cảnh báo tồn kho
    \begin{itemize}
    \item 5.4.1. Cảnh báo sắp hết hàng
    \item 5.4.2. Cảnh báo sắp hết hạn
    \end{itemize}
\end{itemize}

\textbf{6. Quản lý đơn đặt labo}
\begin{itemize}
\item 6.1. Quản lý đối tác labo
    \begin{itemize}
    \item 6.1.1. Thêm đối tác
    \item 6.1.2. Cập nhật thông tin
    \item 6.1.3. Xóa đối tác
    \end{itemize}
\item 6.2. Tạo đơn đặt labo
\item 6.3. Xem danh sách đơn labo
\item 6.4. Cập nhật trạng thái đơn labo
\item 6.5. Hủy đơn labo
\end{itemize}

\textbf{7. Báo cáo và thống kê}
\begin{itemize}
\item 7.1. Báo cáo doanh thu
    \begin{itemize}
    \item 7.1.1. Báo cáo theo tuần
    \item 7.1.2. Báo cáo theo tháng
    \item 7.1.3. Báo cáo theo năm
    \item 7.1.4. Báo cáo tùy chỉnh
    \end{itemize}
\item 7.2. Báo cáo hiệu suất nhân viên
\item 7.3. Báo cáo tài chính
    \begin{itemize}
    \item 7.3.1. Tổng doanh thu
    \item 7.3.2. Tổng chi phí
    \item 7.3.3. Lợi nhuận/Lỗ
    \end{itemize}
\item 7.4. Xuất báo cáo
\end{itemize}

\textbf{8. Quản lý người dùng}
\begin{itemize}
\item 8.1. Đăng ký tài khoản
\item 8.2. Đăng nhập
\item 8.3. Xem thông tin cá nhân
\item 8.4. Cập nhật thông tin cá nhân
\item 8.5. Đổi mật khẩu
\item 8.6. Đăng xuất
\end{itemize}

\vspace{1em}

\textbf{Mô tả chi tiết các chức năng:}

\begin{table}[h]
\centering
\caption{Mô tả các chức năng chính của hệ thống}
\scriptsize
\begin{tabular}{|p{3.5cm}|p{5cm}|p{5cm}|}
\hline
\textbf{Tên chức năng} & \textbf{Mô tả} & \textbf{Đánh giá khả năng thực hiện} \\
\hline
Quản lý phòng khám & Cho phép tạo và quản lý thông tin phòng khám, bao gồm tên, mã code 6 ký tự, quản lý thành viên và thiết lập lương. & Nhân lực: 1 developer. Thời gian: 3 ngày. Công nghệ: Spring Boot + React + MySQL. Môi trường: Web-based, cross-platform. \\
\hline
Thêm bệnh nhân mới & Đăng ký thông tin bệnh nhân bao gồm họ tên, số điện thoại, địa chỉ, ngày sinh, ghi chú. Kiểm tra tính hợp lệ của dữ liệu đầu vào. & Nhân lực: 1 developer. Thời gian: 1 ngày. Công nghệ: Form validation với React, JPA entity validation. Môi trường: Web form. \\
\hline
Tìm kiếm bệnh nhân & Cho phép tìm kiếm bệnh nhân theo tên, số điện thoại hoặc mã ID. Hỗ trợ tìm kiếm không dấu tiếng Việt. & Nhân lực: 1 developer. Thời gian: 1 ngày. Công nghệ: MySQL fulltext search hoặc JPA query. Môi trường: Backend API + Frontend search bar. \\
\hline
Tạo lịch hẹn & Tạo lịch hẹn cho bệnh nhân với bác sĩ cụ thể, kiểm tra xung đột lịch. & Nhân lực: 1 developer. Thời gian: 2 ngày. Công nghệ: DateTime handling, conflict detection algorithm. Môi trường: Web interface với date-time picker. \\
\hline
Xem lịch dạng calendar & Hiển thị lịch hẹn dưới dạng lịch theo ngày/tuần/tháng với khả năng tương tác. & Nhân lực: 1 developer. Thời gian: 3 ngày. Công nghệ: React calendar library (FullCalendar.js). Môi trường: Modern web browsers. \\
\hline
Tạo hồ sơ điều trị & Ghi nhận thông tin điều trị bao gồm chẩn đoán, phương pháp điều trị, tổng chi phí. Liên kết với bệnh nhân và bác sĩ điều trị. & Nhân lực: 1 developer. Thời gian: 2 ngày. Công nghệ: JPA relationships, transaction management. Môi trường: Web form với rich text editor. \\
\hline
Quản lý thanh toán & Ghi nhận các khoản thanh toán, theo dõi công nợ, hỗ trợ thanh toán từng phần. & Nhân lực: 1 developer. Thời gian: 2 ngày. Công nghệ: BigDecimal for currency, payment tracking. Môi trường: Web interface. \\
\hline
Quản lý tồn kho & Quản lý vật tư y tế, thuốc, thiết bị. Theo dõi số lượng tồn, ngày hết hạn, cảnh báo tồn kho thấp. & Nhân lực: 1-2 developers. Thời gian: 5 ngày. Công nghệ: Inventory tracking algorithms, batch management. Môi trường: Web-based inventory system. \\
\hline
Nhập/Xuất kho & Ghi nhận giao dịch nhập/xuất kho, cập nhật số lượng tồn kho tự động. & Nhân lực: 1 developer. Thời gian: 2 ngày. Công nghệ: Database transactions, ACID compliance. Môi trường: Web forms. \\
\hline
Quản lý nhà cung cấp & Lưu trữ thông tin nhà cung cấp bao gồm tên, người liên hệ, điện thoại, email, điều khoản thanh toán. & Nhân lực: 1 developer. Thời gian: 1 ngày. Công nghệ: CRUD operations, JPA. Môi trường: Web interface. \\
\hline
Tạo đơn đặt labo & Tạo đơn đặt hàng cho phòng labo với thông tin mô tả, giá, ngày giao hàng dự kiến. & Nhân lực: 1 developer. Thời gian: 2 ngày. Công nghệ: Order management system, status tracking. Môi trường: Web forms. \\
\hline
Cập nhật trạng thái labo & Theo dõi và cập nhật trạng thái đơn labo (ORDERED, RECEIVED, INSTALLED, CANCELED). & Nhân lực: 1 developer. Thời gian: 1 ngày. Công nghệ: State machine, status workflow. Môi trường: Web interface. \\
\hline
Báo cáo doanh thu & Tính toán và hiển thị doanh thu theo tuần/tháng/năm, so sánh với chi phí, tính lợi nhuận. & Nhân lực: 1 developer. Thời gian: 3 ngày. Công nghệ: SQL aggregation, Recharts for visualization. Môi trường: Dashboard with charts. \\
\hline
Báo cáo hiệu suất nhân viên & Thống kê số lượng điều trị, doanh thu theo từng nhân viên trong khoảng thời gian. & Nhân lực: 1 developer. Thời gian: 2 ngày. Công nghệ: SQL GROUP BY, data aggregation. Môi trường: Web dashboard. \\
\hline
Xác thực và phân quyền & Đăng ký, đăng nhập với JWT, phân quyền Owner/Staff với các chức năng khác nhau. & Nhân lực: 1 developer. Thời gian: 3 ngày. Công nghệ: Spring Security, JWT, role-based access control. Môi trường: Stateless authentication. \\
\hline
\end{tabular}
\end{table}

\subsection{Xây dựng kế hoạch dự án đơn giản}

\textbf{Bảng 1: Phân rã công việc và ước lượng thời gian}

\begin{table}[h]
\centering
\caption{Kế hoạch công việc dự án}
\begin{tabular}{|p{7cm}|p{3cm}|p{2cm}|}
\hline
\textbf{Công việc} & \textbf{Thời gian (giờ)} & \textbf{Số người} \\
\hline
\textbf{1. Thiết lập dự án và cơ sở hạ tầng} & \textbf{32} & \textbf{2} \\
\hline
\quad 1.1. Thiết lập môi trường phát triển & 8 & 1 \\
\hline
\quad 1.2. Thiết kế cơ sở dữ liệu & 12 & 1 \\
\hline
\quad 1.3. Cấu hình Spring Boot backend & 8 & 1 \\
\hline
\quad 1.4. Cấu hình React frontend & 4 & 1 \\
\hline
\textbf{2. Phát triển module xác thực} & \textbf{24} & \textbf{1} \\
\hline
\quad 2.1. Phát triển API đăng ký/đăng nhập & 12 & 1 \\
\hline
\quad 2.2. Tích hợp JWT authentication & 8 & 1 \\
\hline
\quad 2.3. Phát triển UI đăng ký/đăng nhập & 4 & 1 \\
\hline
\textbf{3. Phát triển module quản lý phòng khám} & \textbf{24} & \textbf{1} \\
\hline
\quad 3.1. Backend: CRUD phòng khám & 8 & 1 \\
\hline
\quad 3.2. Backend: Quản lý thành viên & 8 & 1 \\
\hline
\quad 3.3. Frontend: UI quản lý phòng khám & 8 & 1 \\
\hline
\textbf{4. Phát triển module quản lý bệnh nhân} & \textbf{32} & \textbf{2} \\
\hline
\quad 4.1. Backend: CRUD bệnh nhân & 12 & 1 \\
\hline
\quad 4.2. Backend: Tìm kiếm và lọc & 8 & 1 \\
\hline
\quad 4.3. Frontend: Danh sách và form bệnh nhân & 12 & 1 \\
\hline
\textbf{5. Phát triển module quản lý lịch hẹn} & \textbf{40} & \textbf{2} \\
\hline
\quad 5.1. Backend: CRUD lịch hẹn & 12 & 1 \\
\hline
\quad 5.2. Backend: Kiểm tra xung đột lịch & 8 & 1 \\
\hline
\quad 5.3. Frontend: Danh sách lịch hẹn & 8 & 1 \\
\hline
\quad 5.4. Frontend: Calendar view & 12 & 1 \\
\hline
\textbf{6. Phát triển module quản lý điều trị} & \textbf{40} & \textbf{2} \\
\hline
\quad 6.1. Backend: CRUD điều trị & 12 & 1 \\
\hline
\quad 6.2. Backend: Quản lý thanh toán & 12 & 1 \\
\hline
\quad 6.3. Frontend: Form và chi tiết điều trị & 16 & 1 \\
\hline
\textbf{7. Phát triển module quản lý tồn kho} & \textbf{56} & \textbf{2} \\
\hline
\quad 7.1. Backend: Quản lý vật tư & 12 & 1 \\
\hline
\quad 7.2. Backend: Giao dịch nhập/xuất kho & 16 & 1 \\
\hline
\quad 7.3. Backend: Quản lý nhà cung cấp & 8 & 1 \\
\hline
\quad 7.4. Backend: Cảnh báo tồn kho & 8 & 1 \\
\hline
\quad 7.5. Frontend: UI quản lý kho & 12 & 1 \\
\hline
\textbf{8. Phát triển module đơn đặt labo} & \textbf{32} & \textbf{2} \\
\hline
\quad 8.1. Backend: CRUD đơn labo & 12 & 1 \\
\hline
\quad 8.2. Backend: Quản lý đối tác labo & 8 & 1 \\
\hline
\quad 8.3. Frontend: UI quản lý đơn labo & 12 & 1 \\
\hline
\textbf{9. Phát triển module báo cáo thống kê} & \textbf{40} & \textbf{2} \\
\hline
\quad 9.1. Backend: Báo cáo doanh thu & 16 & 1 \\
\hline
\quad 9.2. Backend: Báo cáo hiệu suất nhân viên & 8 & 1 \\
\hline
\quad 9.3. Frontend: Dashboard và biểu đồ & 16 & 1 \\
\hline
\textbf{10. Kiểm thử và hoàn thiện} & \textbf{48} & \textbf{3} \\
\hline
\quad 10.1. Kiểm thử đơn vị (Unit testing) & 16 & 1 \\
\hline
\quad 10.2. Kiểm thử tích hợp & 16 & 2 \\
\hline
\quad 10.3. Sửa lỗi và tối ưu hóa & 16 & 2 \\
\hline
\textbf{11. Triển khai và tài liệu} & \textbf{24} & \textbf{2} \\
\hline
\quad 11.1. Viết tài liệu hướng dẫn sử dụng & 8 & 1 \\
\hline
\quad 11.2. Viết tài liệu kỹ thuật & 8 & 1 \\
\hline
\quad 11.3. Chuẩn bị triển khai & 8 & 2 \\
\hline
\hline
\textbf{TỔNG CỘNG} & \textbf{392} & \textbf{-} \\
\hline
\end{tabular}
\end{table}

\vspace{1em}

\textbf{Bảng 2: Quản lý rủi ro dự án}

\begin{table}[h]
\centering
\caption{Bảng quản lý rủi ro dự án}
\scriptsize
\begin{tabular}{|p{2.5cm}|p{2.5cm}|p{2cm}|p{1.8cm}|p{1.8cm}|p{2.5cm}|}
\hline
\textbf{Công việc/Hoạt động} & \multicolumn{3}{c|}{\textbf{Xác định rủi ro}} & \multicolumn{2}{c|}{\textbf{Quản lý rủi ro}} \\
\cline{2-6}
 & \textbf{Mối nguy} & \textbf{Rủi ro} & \textbf{Mức độ} & \textbf{Chiến lược} & \textbf{Biện pháp} \\
\hline
Thiết kế cơ sở dữ liệu & Thiết kế không đầy đủ, thiếu ràng buộc & Phải thiết kế lại sau này, ảnh hưởng toàn bộ dự án & Cao & Tránh & Đầu tư thời gian phân tích kỹ, review bởi senior developer, tham khảo best practices \\
\hline
Phát triển backend API & Thiếu kinh nghiệm với Spring Security & API không an toàn, dễ bị tấn công & Cao & Giảm thiểu & Học và áp dụng Spring Security best practices, sử dụng JWT, code review thường xuyên \\
\hline
Tích hợp frontend-backend & Sai lệch về cấu trúc dữ liệu giữa FE và BE & Lỗi khi gọi API, mất thời gian debug & Trung bình & Giảm thiểu & Định nghĩa rõ ràng DTO, sử dụng TypeScript, viết API documentation chi tiết \\
\hline
Quản lý thanh toán & Sai sót trong tính toán số tiền, công nợ & Sai sót tài chính, mất uy tín & Cao & Tránh & Sử dụng BigDecimal, kiểm thử kỹ các case thanh toán, review code cẩn thận \\
\hline
Quản lý tồn kho & Race condition khi nhiều người xuất kho cùng lúc & Số liệu tồn kho không chính xác & Cao & Giảm thiểu & Sử dụng database transaction, optimistic/pessimistic locking \\
\hline
Hiển thị lịch hẹn & Xung đột múi giờ, format ngày tháng & Hiển thị sai thời gian hẹn & Trung bình & Giảm thiểu & Chuẩn hóa timezone UTC, sử dụng LocalDateTime, test kỹ với nhiều timezone \\
\hline
Báo cáo thống kê & Query phức tạp, chậm với dữ liệu lớn & Hiệu suất kém, người dùng phải chờ lâu & Trung bình & Giảm thiểu & Tối ưu query SQL, tạo index, cache kết quả, sử dụng pagination \\
\hline
Tìm kiếm tiếng Việt & Không tìm được khi nhập không dấu & Trải nghiệm người dùng kém & Thấp & Chấp nhận & Document hướng dẫn người dùng, hoặc thêm chức năng convert không dấu nếu có thời gian \\
\hline
Kiểm thử & Thiếu thời gian kiểm thử đầy đủ & Nhiều lỗi khi đưa vào sử dụng & Cao & Giảm thiểu & Lập kế hoạch kiểm thử từ đầu, viết unit test song song với code, automated testing \\
\hline
Triển khai & Server không ổn định, thiếu kiến thức DevOps & Ứng dụng không chạy được trên production & Trung bình & Giảm thiểu & Sử dụng cloud platform đáng tin cậy (AWS, Heroku), docker containerization, document deployment process \\
\hline
Yêu cầu thay đổi & Khách hàng thay đổi yêu cầu giữa chừng & Mất thời gian, vượt budget & Trung bình & Giảm thiểu & Xác định rõ yêu cầu từ đầu, có hợp đồng rõ ràng, sử dụng Agile để linh hoạt \\
\hline
Quản lý thời gian & Ước lượng sai thời gian, trễ deadline & Chất lượng kém, khách hàng không hài lòng & Trung bình & Giảm thiểu & Ước lượng thêm 20-30\% buffer time, theo dõi tiến độ hàng tuần, ưu tiên các chức năng core \\
\hline
Thành viên nghỉ việc & Developer chính rời dự án giữa chừng & Mất kiến thức, chậm tiến độ & Cao & Giảm thiểu & Document code tốt, pair programming, knowledge sharing sessions, backup developer \\
\hline
\end{tabular}
\end{table}
